\cotefile{93}
\scanpage{207}

On a donc

\diagram{\begin{tikzcd}[ampersand replacement=\&,column sep=small,row sep=small]
\llap{(2.1)\qquad}0 \arrow[r] \& P^{ci}_{X/S}\arrow[r] \& P^c_{X/S}\arrow[r,"\text{courbure}"] \& f_*^{\mathrm{fl}}(Z^2_{X/S}) \\
\&\&\& f_*^{\mathrm{fl}}(\Omega^2_{X/S})\arrow[u,phantom,"\rotatebox{-90}{$\subset$}" description]
\end{tikzcd}}

ou, si on préfère, dans le cas où $X$ est lisse sur $S$.

\[
0\longrightarrow\frac{f_*^{\mathrm{fl}}(\underline{Z}^1_{X/S})}{\underline{\mathrm{Im}}\,f_*^{\mathrm{fl}}(\underline O^*_{X/S})}
\longrightarrow P^{ci}_{X/S}\longrightarrow\underline{\mathrm{Pic}}_{X/S}
\longrightarrow R^{2}f_*^{\mathrm{fl}}(0\to\underline\Omega^1_{X/S}\to\underline\Omega^2_{X/S}\to\unsure{\cdots})
\tag*{\unsure{(1.2)|(2.2)}}
\]

[Dans le cas non lisse, on remplace le dernier terme par \deleted{$R^{\unsure{2}}f_*^{\mathrm{cris}}(\unsure{\mathscr J|\mathscr I})$} le faisceau associé à $S'\longmapsto R^2 f_{\underline{S'}\,\mathrm{com}}(\mathscr J_{X_{S'}/S'})$ \inserted{idéal d'augmentation du $(X_{S'}/S')\,\mathrm{cris}$}]

% slanted
\margin{\unsure{pourrait|peut entrer|pour entrer|remis} dans $1^\circ$)}

\uline{Cas particulier} Plaçons nous \unsure{sur|dans} le cas où $\underline{\mathrm{Pic}}^0_{X/S}$ est un schéma abélien $B$, dont le dual est noté $A^0=\underline{\mathrm{Alb}}^0_{X/S}$.

On a alors $X\to\underline{\mathrm{Alb}}^1_{X/S}=A^1$ (torseur sous $A^0$), et supposons que
\diagram{\begin{tikzcd}[ampersand replacement=\&,column sep=tiny,row sep=tiny]
X\arrow[rr,dash]\arrow[dr,dash,"f"'] \& \& A^1\arrow[dl,dash,"\unsure{g^1|g^2}"] \\
\& S \&
\end{tikzcd}}
induise un \uline{isom} $f_*^{\mathrm{fl}}(\Omega^1_{X/S})\longleftarrow {g^1}_*^{\mathrm{fl}}(\Omega^1_{A^1/S})$.

Il suffit pour ceci que \inserted{$f$ pr. plat, que} $\Omega^1_{X/S}$ \deleted{soit} \inserted{plat et} coh.\ plat en dim 0, et que l'on ait isom des $H^0(A^1_s,\Omega^1_{A^1_s/s})\to H^0(X,\Omega^1_{X_s/s})$ fibre à fibre. C'est le cas en particulier si $X$ est un torseur sous schéma abélien (bien sûr), ou si $X$ est une courbe relative propre et lisse.

Alors $P^{c\unsure{0|o}}_{X/S}$ \inserted{de (1.3)} est une extension du schéma abélien $B$ par un fibré vectoriel, et on constate que c'est l'\uline{extension universelle} [compte tenu que $B^*\simeq A$, \deleted{${\unsure{\underline{\mathrm{Lie}}}}_{B^*/\unsure{S}}$} $\underline{\mathrm{Lie}}(B^*)^\vee=\underline{\mathrm{Lie}}(A)^\vee\simeq {g^0}_*(\Omega^1_{\unsure{A^1|A}/S})\simeq {g^1}_*(\Omega^1_{A^1/S})\simeq f_*(\Omega^1_{A^1/S})\ldots$

% archive sheet number: 206
